%%
%% Progress Report
%% Multimodal Distillation Method for Heart Disease Classification
%%
\documentclass[sigconf,nonacm]{acmart}

\usepackage{booktabs}
\usepackage{graphicx}
\usepackage{xcolor}
\usepackage{enumitem}
\usepackage{multirow}
\usepackage{array}
\usepackage{url}

\settopmatter{printacmref=false}
\renewcommand\footnotetextcopyrightpermission[1]{}
\pagestyle{plain}

\title{Multimodal Distillation Method for Heart Disease Classification}
\subtitle{Progress Update}

\author{Curtis Choy}
\affiliation{\institution{Texas A\&M University}}
\email{cthcurtis@tamu.edu}

\author{Dhanush Shekar}
\affiliation{\institution{Texas A\&M University}}
\email{dshekar@tamu.edu}

\author{Alexandra Tabacaru}
\affiliation{\institution{Texas A\&M University}}
\email{alex.tabac@tamu.edu}

\begin{document}
\maketitle

%% ──────────────────────────────────────────────────────────────
\begin{abstract}
We aim to make the ECG — a fast, inexpensive first-tier diagnostic signal
— as informative as a full multimodal assessment combining ECG, chest
X-ray (CXR), and clinical text. Using knowledge distillation, a pretrained
multimodal teacher (MoRE~\cite{thapa2024more}) transfers its diagnostic
knowledge to a lightweight single-lead ECG student (MobileNetV3-Small,
{\raise.17ex\hbox{$\scriptstyle\sim$}}2.54M parameters). This report
describes progress through Phase 4: data procurement from PhysioNet,
manifest construction and partial download under HPRC inode constraints,
MoRE framework adaptation, ECG rhythm label extraction, and teacher
embedding caching. The student model architecture has been designed and
implemented with dual classification heads — a primary ECG rhythm
classification head (10 classes) and a secondary pulmonary risk indicator
head (4 CheXpert classes) — and distillation training has been initiated.

\subsection*{Availability}
Code, preprocessing scripts, SLURM job definitions, and configuration
files are available on GitHub at
\url{https://github.com/Dhanush006/more-clinical-distill} (branch:
\texttt{more-KD}). The processed dataset manifests and embedding caches
are not redistributable (MIMIC data access requires PhysioNet credentials);
however, the scripts to reproduce all preprocessing steps from a credentialed
PhysioNet download are fully documented in the repository \texttt{README}.
\end{abstract}

%% ──────────────────────────────────────────────────────────────
\section{Introduction}

Heart disease remains a leading cause of global mortality, accounting for
roughly 19.8 million deaths annually~\cite{who_cvd}. Missed diagnostic
opportunities — often attributable to limited specialist access or
single-modality assessment — contribute to delayed treatment and
preventable mortality~\cite{kwok2022}. Multimodal AI systems that fuse
ECG, CXR, and clinical text have demonstrated improved diagnostic
accuracy~\cite{thapa2024more}, but their computational footprint limits
deployment in resource-constrained or edge settings.

Knowledge distillation~\cite{hinton2015distilling} offers a path forward:
a capable multimodal teacher imparts its learned representations to a
lightweight student that operates on a single modality at inference time.
Our project applies this paradigm to cardiovascular diagnosis, distilling
the MoRE multimodal teacher into a MobileNetV3-Small student that uses
only single-lead ECG input.

\section{Related Work}

\subsection{Unimodal AI for Cardiovascular Diagnosis}
CNN and transformer-based models have been widely applied to ECG
arrhythmia and cardiac abnormality detection~\cite{avanzato2020,noor2025}.
These approaches benefit from simplicity but cannot leverage richer
multimodal representations.

\subsection{Multimodal Medical AI}
The MoRE framework~\cite{thapa2024more} aligns ECG, CXR, and clinical text
embeddings using InfoNCE contrastive learning across a shared 128-dim
embedding space. It demonstrates that multimodal alignment yields stronger
representations than any single modality alone.

\subsection{Knowledge Distillation in Medical AI}
Hossain et al.~\cite{hossain2025} applied multimodal distillation to tumor
segmentation, aggregating knowledge from multiple modality-specific teacher
models into a single student with fewer parameters. Application of this
paradigm to cardiovascular disease with ECG signals remains unexplored.

%% ──────────────────────────────────────────────────────────────
\section{Proposed Method}

\subsection{Multimodal Teacher}
The MoRE teacher encodes three modalities:
\begin{itemize}[noitemsep]
  \item \textbf{ECG encoder:} ViT-Base with custom Conv1D patch embedding
        (12-lead, $12\times1000$ input at 100\,Hz).
  \item \textbf{CXR encoder:} ViT-Base with DropKey attention.
  \item \textbf{Text encoder:} RoBERTa-base-PM-M3 with LoRA ($r=16$,
        0.6\% of parameters trainable).
\end{itemize}
Modality projectors map each encoder to a shared 128-dim space.
Pre-training uses InfoNCE contrastive loss across ECG, CXR, and text
pairs. The teacher ($\sim$280M parameters) is fully frozen during student training.

\subsection{Student Model}
The student (Figure~\ref{fig:arch}) is a MobileNetV3-Small backbone
adapted for 1-D single-lead ECG input:

\begin{enumerate}[noitemsep]
  \item \textbf{Channel adapter:} Conv1d $1\to16\to3$, BN, ReLU — maps
        single-lead time series to pseudo-3-channel representation.
  \item \textbf{Backbone:} MobileNetV3-Small (\texttt{timm} library),
        producing a 576-dim feature vector.
  \item \textbf{Projector:} Linear $576\to128$, matching teacher embedding
        dimensionality.
  \item \textbf{ECG head (primary):} Linear $128\to10$ — 10 ECG rhythm
        classes.
  \item \textbf{Pulm head (secondary):} Linear $128\to4$ — 4 CheXpert
        pulmonary classes, learned indirectly via embedding alignment with
        the multimodal teacher.
\end{enumerate}
Total student parameters: $\sim$2.54M — a 110$\times$ compression of the teacher.

\subsection{Distillation Objectives}
Training minimizes a weighted combination:
\begin{equation}
  \mathcal{L} = \alpha_{\text{ecg}}\,\mathcal{L}_{\text{ecg}}^{\text{BCE}}
              + \alpha_{\text{pulm}}\,\mathcal{L}_{\text{pulm}}^{\text{BCE}}
              + \gamma\,\mathcal{L}_{\text{align}}
\end{equation}
where $\mathcal{L}_{\text{align}} = 1 - \overline{\cos}(\mathbf{z}_s, \mathbf{z}_t)$
is the cosine embedding alignment loss. Default weights:
$\alpha_{\text{ecg}}=1.0$, $\alpha_{\text{pulm}}=0.5$, $\gamma=0.5$.
Early stopping is on the primary ECG macro AUC (patience 15 epochs).

%% ──────────────────────────────────────────────────────────────
\section{Progress and Current Results}

\subsection{Phase 0 — Repository Audit and Planning}

A full audit of the MoRE codebase identified several integration barriers
that required fixing before training could proceed:
\begin{itemize}[noitemsep]
  \item Missing imports in \texttt{pretrain\_multimodel.py} (\texttt{torch},
        \texttt{tqdm}, AMP utilities).
  \item Early-\texttt{return} bug in \texttt{preprocess\_notes.py}
        \texttt{get\_clinical\_xray()} causing only one radiology report
        to be processed instead of iterating all records.
  \item All data paths hardcoded as relative strings throughout the
        preprocessing scripts — incompatible with HPC directory layout.
  \item RoBERTa model path hardcoded in \texttt{utils/build\_model.py}.
\end{itemize}
The 49,076-row manifest (\texttt{matched\_patients\_v1.csv}) was audited:
zero null values, 87.2\% same-day ECG–CXR pairs, unique per-subject.
An implementation plan, assumptions log, and full SLURM job chain were
designed prior to writing any training code.

\subsection{Phase 1 — Environment and Reproducibility}

A reproducible HPC environment was established on the Texas A\&M Grace
cluster (A100 40\,GB GPUs). All configurable paths were centralized into
\texttt{configs/paths.yaml}. Credentials for PhysioNet download are
managed via a \texttt{.env} file (never committed). The full project
directory structure — including \texttt{data/}, \texttt{manifests/},
\texttt{scripts/}, \texttt{slurm/}, \texttt{configs/}, \texttt{logs/},
\texttt{outputs/}, and \texttt{distill/} — was created and documented.

\subsection{Phase 2 — Data Procurement (Challenges)}

\textbf{Manifest-driven download.}
\texttt{scripts/build\_download\_lists.py} generated:
98,152 ECG file paths (\texttt{.hea} + \texttt{.dat} per record) and
49,076 CXR JPEG paths, plus all required metadata CSVs
(\texttt{mimic-cxr-2.0.0-chexpert.csv}, \texttt{mimic-cxr-2.0.0-split.csv},
\texttt{machine\_measurements.csv}, etc.).

\textbf{HPRC inode quota constraint.}
The Grace cluster enforces a per-user inode limit ($\sim$1M inodes).
The combined MIMIC-IV-ECG and MIMIC-CXR-JPG directory tree — including
radiology report \texttt{.txt} files, metadata CSVs, and signal binaries
— exceeds this limit. The download SLURM job was terminated by quota
enforcement before CXR JPEG files or radiology report text files could be
stored. Table~\ref{tab:download} summarizes the outcome.

\begin{table}[h]
\caption{Download outcome for 49,076-pair manifest}
\label{tab:download}
\begin{tabular}{lrr}
\toprule
File type & Downloaded & Missing \\
\midrule
ECG metadata CSVs & 2 / 2 & 0 \\
CXR metadata CSVs & 3 / 3 & 0 \\
ECG signal files (.hea+.dat) & 28,745 & 20,331 \\
CXR JPEG images & 0 & 49,076 \\
Radiology report .txt & 0 & $\sim$49,000 \\
\bottomrule
\end{tabular}
\end{table}

\textbf{Mitigation strategy.}
Because the student model uses only ECG signals at inference time, and
because the teacher embedding cache requires only the teacher's ECG
encoder (not CXR), the 28,745 downloaded ECG records are sufficient to
proceed with the distillation pipeline. All 49,076 records remain in
the \texttt{.npy} preprocessing output (paths stored); records with
missing ECG files are filtered at dataset load time.

For xray notes (used only in teacher pretraining, not distillation),
MoRE's built-in CheXpert-label-to-text fallback
(\texttt{generate\_xray\_note()}) is applied to all records.

\subsection{Phase 3 — MoRE Compatibility and Preprocessing}

\textbf{Bug fixes applied to MoRE:}
All hardcoded paths were made configurable via \texttt{configs/paths.yaml}.
The \texttt{preprocess\_notes.py} early-return bug was corrected.
Missing imports in \texttt{pretrain\_multimodel.py} were added.

\textbf{Item format extension.}
The original MoRE 6-element item format was extended to 7 elements:
\begin{center}
\texttt{[xray\_path, ecg\_stem, xray\_note, ecg\_note, ecg\_labels(10,), pulm\_labels(4,), split]}
\end{center}

\textbf{ECG rhythm label extraction.}
The \texttt{machine\_measurements.csv} report columns
(\texttt{report\_0}–\texttt{report\_6}) were parsed using compiled
case-insensitive regular expressions to extract 10-class multi-label
binary rhythm vectors. The ``Normal ECG'' pattern uses a negative
lookbehind (\texttt{(?<![a-z])normal ecg\textbackslash b}) to avoid
matching ``Abnormal ECG.'' Table~\ref{tab:ecg_class} shows the class
distribution across all 49,076 records.

\begin{table}[h]
\caption{ECG rhythm label distribution across all splits}
\label{tab:ecg_class}
\begin{tabular}{lrrr}
\toprule
Class & Total & \% of records \\
\midrule
Normal & 11,053 & 22.5\% \\
Sinus tachycardia & 6,720 & 13.7\% \\
Sinus bradycardia & 4,456 & 9.1\% \\
Atrial fibrillation & 3,978 & 8.1\% \\
LVH & 3,863 & 7.9\% \\
RBBB & 3,345 & 6.8\% \\
ST ischemia & 2,847 & 5.8\% \\
AV block & 2,627 & 5.4\% \\
LBBB & 1,055 & 2.1\% \\
ST elevation MI & 480 & 1.0\% \\
\midrule
No label assigned & 17,166 & 35.0\% \\
\bottomrule
\end{tabular}
\end{table}

The dataset exhibits a notable class imbalance: \textit{ST elevation MI}
(480 records) and \textit{LBBB} (1,055 records) are severely
underrepresented. Class-weighted BCE is planned as an ablation.
35\% of records carry no ECG rhythm label; these contribute to the
pulmonary head and embedding alignment loss but not the ECG BCE term.

\textbf{Preprocessing output:}
\begin{center}
\begin{tabular}{lr}
Split & Records \\
\hline
Train & 48,423 \\
Val & 379 \\
Test & 274 \\
\textbf{Total} & \textbf{49,076} \\
\end{tabular}
\end{center}

\subsection{Phase 4 — Teacher ECG Embedding Cache}

\texttt{distill/cache\_teacher\_embeddings.py} runs the frozen MoRE
teacher's ECG encoder and 128-dim projector over all available ECG
records in batch mode, saving results to
\texttt{data/processed/teacher\_ecg\_embeddings.npy}
($28{,}745 \times 128$, float32, $\sim$14\,MB).

A study-ID index (\texttt{teacher\_ecg\_study\_ids.npy}) enables
efficient lookup during distillation dataset construction.

\textbf{NaN embedding issue.}
502 of the 28,745 cached embeddings contained NaN values, caused by
ViT encoder instability on a small subset of ECG records (likely
unusual signal morphologies). This would have propagated into the
embedding alignment loss and produced NaN training losses on the very
first backward pass. The fix is applied in two layers:
\begin{enumerate}[noitemsep]
  \item \textbf{Dataset load:} \texttt{np.nan\_to\_num(raw, nan=0.0)} replaces all NaN embeddings with the zero vector.
  \item \textbf{Training loop:} A NaN guard (\texttt{if not torch.isfinite(loss): skip batch}) prevents NaN gradient propagation.
\end{enumerate}
Gradient clipping (\texttt{clip\_grad\_norm\_, max\_norm=1.0}) is also
applied as an additional stability measure.

\subsection{Phase 5 — Student Model Design (In Progress)}

The dual-head MobileNetV3-Small student has been fully implemented
(\texttt{distill/student\_model.py}, \texttt{distill/distill\_dataset.py},
\texttt{distill/distill\_train.py}, \texttt{distill/evaluate\_student.py}).
The distillation training SLURM job has been submitted to the Grace
GPU partition. Full training results and ablation studies will be
reported in the final submission.

%% ──────────────────────────────────────────────────────────────
\section{Barriers and New Ideas}

\subsection{Barriers Encountered}

\textbf{HPRC inode quota.} The most significant infrastructure barrier.
With only 28,745 ECG records available and no CXR files, the training
set is approximately 59\% of the intended size. Potential resolution
paths include: (a) requesting a quota increase from HPRC, (b) downloading
only the specific CXR files needed (already have exact paths), or
(c) using group-shared scratch storage.

\textbf{NaN teacher embeddings.} A non-obvious signal quality issue
that would have silently caused all training loss values to be NaN from
epoch 1. Required numpy-level inspection of the cached embedding array
to diagnose.

\textbf{SLURM dependency chaining.} Multi-line \texttt{sbatch} output
broke naive variable capture for job chaining. Fixed by piping through
\texttt{grep "Submitted batch job" | awk '\{print \$NF\}'}.

\textbf{timm API incompatibility.} The MoRE teacher uses older
\texttt{timm} API conventions for ViT model creation that were
deprecated in newer versions. Compatibility shims were added.

\subsection{New Ideas for Ablations}

\textbf{Class-weighted BCE.} ST elevation MI (480 records, 1.0\%) and
LBBB (1,055 records, 2.1\%) are likely to be underfit with uniform BCE.
Class weights inversely proportional to frequency will be tested.

\textbf{Multi-lead student.} Instead of a single lead, concatenate
2–3 leads as input channels. This keeps the student lightweight but
may substantially improve rhythm classification AUC on complex classes
like AF and AV block.

\textbf{Soft label distillation.}
The current implementation uses hard labels for both ECG and pulm heads.
Adding KL divergence against teacher soft logits (temperature-scaled
softmax) as proposed in the original plan~\cite{hinton2015distilling}
may improve calibration and secondary head performance.

%% ──────────────────────────────────────────────────────────────
\section{Action Items and Expected Outcomes}

\begin{table}[h]
\caption{Remaining action items and expected outcomes}
\label{tab:action}
\begin{tabular}{p{3.5cm}p{4cm}}
\toprule
Action item & Expected outcome \\
\midrule
Complete distillation run (Lead I baseline) &
  Per-class ECG AUC-ROC on test set; embedding cosine similarity to teacher \\
\addlinespace
Lead ablation (Lead I vs II vs V2) &
  Identify strongest single lead for ECG rhythm classification \\
\addlinespace
Loss weight ablation ($\gamma$, $\alpha_\text{pulm}$) &
  Quantify contribution of embedding alignment and pulm head to ECG performance \\
\addlinespace
Class-weighted BCE &
  Improve AUC for ST elevation MI and LBBB \\
\addlinespace
Teacher linear probe benchmark &
  Establish ceiling AUC from teacher embeddings alone (LogReg on cached embeddings) \\
\addlinespace
Parameter count \& latency measurement &
  Demonstrate 110$\times$ compression and $<$5\,ms inference on a single ECG record \\
\addlinespace
Final report + presentation &
  April 23–30: write-up with full ablation table, error analysis, conclusion \\
\bottomrule
\end{tabular}
\end{table}

%% ──────────────────────────────────────────────────────────────
\section{Updated Project Timeline}

Figure~\ref{fig:gantt} shows the updated Gantt chart reflecting actual
execution. The data procurement phase extended beyond the original schedule
due to the HPRC inode quota issue, which required iterative debugging
and a revised download strategy. Teacher embedding caching (a new Phase 4
task not in the original proposal) added $\sim$1 week. Student model
design, implementation, and initial training are on schedule relative to
the revised timeline.

\begin{figure}[h]
\centering
\fbox{\parbox{0.9\columnwidth}{
\small
\textbf{Phase 1 — Planning \& Proposal} \hfill Feb 11 – Mar 2 \checkmark \\[2pt]
\textbf{Phase 2 — Data \& Teacher Groundwork} \hfill Mar 9 – Apr 4 \\
\quad 2.1 Data access + manifest audit \hfill \checkmark \\
\quad 2.2 Download pipeline (partial, inode limit) \hfill \checkmark \\
\quad 2.3 MoRE preprocessing + ECG label parsing \hfill \checkmark \\
\quad 2.4 Teacher ECG embedding cache \hfill \checkmark \\[2pt]
\textbf{Phase 3 — Student Model \& Distillation} \hfill Apr 5 – Apr 22 \\
\quad 3.1 Student design \& implementation \hfill \checkmark \\
\quad 3.2 Distillation training (Lead I baseline) \hfill In progress \\
\quad 3.3 Lead ablation study \hfill Pending \\
\quad 3.4 Loss ablation study \hfill Pending \\[2pt]
\textbf{Phase 4 — Results, Polishing, Presentation} \hfill Apr 23 – Apr 30 \\
\quad 4.1 Final evaluation \& error analysis \hfill Pending \\
\quad 4.2 Write-up \hfill Pending \\
\quad 4.3 Presentation \hfill Pending \\
}}
\caption{Updated project timeline (as of 2026-03-30)}
\label{fig:gantt}
\end{figure}

%% ──────────────────────────────────────────────────────────────
\section{Individual Contributions}

\textbf{Curtis Choy} led data access and preprocessing infrastructure:
designed and implemented the manifest-driven download pipeline
(\texttt{scripts/build\_download\_lists.py},
\texttt{slurm/download\_subset.slurm}), HPRC environment configuration,
and the layout verification script. Conducted initial manifest inspection
and data quality audit.

\textbf{Dhanush Shekar} led modeling and training infrastructure:
adapted the MoRE framework (bug fixes, path refactoring), implemented
ECG rhythm label extraction, built the full distillation codebase
(\texttt{distill/} module), designed the dual-head student architecture,
authored SLURM job chains for teacher caching and distillation training,
and diagnosed/fixed the NaN embedding issue.

\textbf{Alexandra Tabacaru} led evaluation design and documentation:
designed the teacher linear probe benchmark methodology, managed the
project timeline and updated Gantt chart, wrote and maintained
\texttt{docs/} documentation, and led early-stopping strategy and
evaluation metric selection (macro AUC-ROC as primary metric).

%% ──────────────────────────────────────────────────────────────
\bibliographystyle{ACM-Reference-Format}
\begin{thebibliography}{9}

\bibitem{who_cvd}
World Health Organization.
\newblock Cardiovascular diseases (CVDs).
\newblock \url{https://www.who.int/news-room/fact-sheets/detail/cardiovascular-diseases-(cvds)}, 2023.

\bibitem{kwok2022}
Chun Shing Kwok et al.
\newblock Missed opportunities in the diagnosis of heart failure: Evaluation of pathways to determine sources of delay to specialist evaluation.
\newblock \textit{Current Heart Failure Reports}, 19(4):247–253, 2022.

\bibitem{naidoo2026}
Vasanthrie Naidoo, Lavanya Madamshetty, and Suresh Babu Krishna.
\newblock Artificial intelligence for cardiology: From diagnosis to management.
\newblock \textit{Therapeutic Advances in Cardiovascular Disease}, 20, 2026.

\bibitem{thapa2024more}
Samrajya Thapa, Koushik Howlader, Subhankar Bhattacharjee, and Wei le.
\newblock MoRE: Multi-modal contrastive pre-training with transformers on X-rays, ECGs, and diagnostic report.
\newblock \textit{arXiv preprint arXiv:2410.16239}, 2024.

\bibitem{hossain2025}
Khondker Fariha Hossain, Sharif Amit Kamran, Joshua Ong, and Alireza Tavakkoli.
\newblock Enhancing efficient deep learning models with multimodal, multi-teacher insights for medical image segmentation.
\newblock \textit{Scientific Reports}, 15(1), 2025.

\bibitem{avanzato2020}
Avanzato R and Beritelli F.
\newblock Automatic ECG diagnosis using convolutional neural network.
\newblock \textit{Electronics}, 9(6):951, 2020.

\bibitem{noor2025}
Noor N, Bilal M, Abbasi SF, Pournik O, and Arvanitis TN.
\newblock A novel transformer-based approach for cardiovascular disease detection.
\newblock \textit{Frontiers in Digital Health}, 7:1548448, 2025.

\bibitem{hinton2015distilling}
Geoffrey Hinton, Oriol Vinyals, and Jeff Dean.
\newblock Distilling the knowledge in a neural network.
\newblock \textit{arXiv preprint arXiv:1503.02531}, 2015.

\end{thebibliography}

\end{document}
